\section{Implementation details}
\label{sec:details}

\subsection{The event guards}
\label{sec:guards}

Slopes are tested against a floor of $\sqrt{\varepsilon_{\text{mach}}}$ in the
problem's own units, not against the user tolerance, which would let a tiny but genuine
slope go unnoticed. With $m = \lVert\mu\rVert_\infty$ and
$\sigma = \lVert\Sig\mathbf{1}\rVert_\infty/n$, $\lambda$ has the scale $\sigma/m$.
Weight slopes $\beta_i$ and slack slopes $g_j\trans\beta$ carry units of $1/\lambda$
and are tested against $\sqrt{\varepsilon_{\text{mach}}}\,m/\sigma$; multiplier slopes
$\delta_i$, $\eta^{\beta}_j$ carry the units of $\mu$ and are tested against
$\sqrt{\varepsilon_{\text{mach}}}\,m$. The $\lambda$ window is a relative $10^{-10}$
taken against $\max(\lvert\lambda\rvert, \sigma/m)$, so that it does not vanish as
$\lambda \to 0$.

\label{sec:invariance}\label{sec:moreexp}%
The trace does not depend on the units (\target{validate-scaling}). Rescaling the returns
by $c$ and the covariance by $s$, each from $10^{-6}$ to $10^{6}$, leaves all $29$
turning points of the $40$-asset problem in place: the weights agree to
$8\times10^{-15}$, and $\lambda c/s$ to a relative $1.2\times10^{-13}$. Multiplying the
slope floors by any factor from $10^{-6}$ to $10^{4}$ leaves the traces of the $40$- and
$494$-asset S\&P 500 problems, the $600$-day window and the $320$-asset factor market bit
for bit unchanged. Only at $10^{6}$ does the floor start to swallow genuine slopes: the
$40$-asset trace is still unchanged, the $600$-day window's weights move by
$7\times10^{-2}$, and the trace of the $494$-asset problem and of the factor market is
declined. So on these inputs genuine slopes and round-off are about ten orders of
magnitude apart. A turning point
lost to the floor would leave the later ones feasible but not optimal, which the KKT
certificate reports as a stationarity or dual-feasibility residual.

\subsection{The gross-exposure lift}
\label{sec:liftdetail}

Both legs of an asset are never free together on the path. Write $g_i$ for the $i$-th
entry of $\Sig w - \lambda\mu + A\trans\nu + G\trans\eta_G$ and $\eta\ge 0$ for the
multiplier of the cap; the reduced gradients of the two legs are $g_i + \eta$ and
$-g_i + \eta$. A free long leg has $g_i = -\eta$, so its short leg's reduced gradient
is $2\eta$; a long leg at its upper bound has $g_i + \eta \le 0$, so its short leg's is
at least $2\eta$ (and symmetrically). The idle leg can therefore enter only where $\eta$
vanishes, at the release of the cap, and the event search drops that competing leg
event so the release fires.

At the first vertex, overlapping legs (the cap tight with a zero multiplier) are removed
by subtracting $\min(x_i^+, x_i^-)$ from both, which leaves $w$ and every other
constraint unchanged; a cap at the smallest feasible gross exposure ($c = 1$ under a
full budget), which fails Slater's condition, is handled by solving
$\min \mathbf{1}\trans x$ first and pinning every leg with a positive reduced cost.

\subsection{The degenerate first vertex}
\label{sec:vertexdetail}

The maximum-return vertex can be degenerate: every weight on a bound (any inequality row on a
long-only problem whose best asset takes the whole budget), or a duplicated tight row.
Its partition is then read off the linear program's duals, generalising the budget rule
of \S\ref{sec:first}. At a degenerate vertex these are not unique, and the solution the simplex
reports can show too few zero reduced costs to free enough assets. For example, with
three assets, $\mu = (3, 2, 1)$, long-only, the budget and the row $w_2 = 0$, HiGHS
reports duals that free only the first asset at $w = (1, 0, 0)$, whereas the equally
optimal multipliers $1$ on both rows free the other two. The code therefore first moves
the reported duals to a vertex of the optimal dual set, stepping along directions that
keep the tight dual constraints tight until they determine the multipliers. A basis is
read off that dual vertex: the assets strictly inside their bounds, then the tight rows
with zero multiplier, which are released, then the assets on a bound with zero reduced
cost, in order of decreasing return, which are freed. The active rows then have a square,
nonsingular free block, the first segment starts at the vertex itself, and the implied
multipliers are those of the dual vertex, so the partition is dual feasible. The optimal
dual set has no vertex, and the vertex is refused, when a tight row touches only assets
fixed at both bounds.

Nearly dependent equality rows, such as $\mathbf{1}\trans w = 1$ and
$(\mathbf{1} + \epsilon v)\trans w = 1$, which encode $v\trans w = 0$ only at scale
$\epsilon$, would be enforced by the linear program only to its tolerance divided by
$\epsilon$. When the row-normalised equality matrix has a reciprocal condition number
below $10^{-6}$ they are handed to HiGHS in orthonormal rows, $Q\trans w = R^{-\mathsf{T}}
b$ with $A\trans = QR$.
